% Report: floating-strike lookback options, closed form and Monte Carlo engines.
% Build: pdflatex report.tex (twice). Figures read their data from data/*.dat, produced by the page's JavaScript engine.
\documentclass[11pt,a4paper]{article}
\usepackage[a4paper,margin=2.4cm]{geometry}
\usepackage[T1]{fontenc}
\usepackage[utf8]{inputenc}
\usepackage{lmodern}
\usepackage{microtype}
\usepackage{amsmath,amssymb,amsthm,mathtools,bm}
\usepackage{booktabs}
\usepackage{enumitem}
\usepackage{xcolor}
\usepackage{pgfplots}
\pgfplotsset{compat=1.17}
\usepackage{caption}
\usepackage[colorlinks=true,linkcolor=blue!50!black,citecolor=blue!50!black,urlcolor=blue!50!black]{hyperref}
\usepackage{cleveref}
\hypersetup{pdftitle={Floating-strike lookback options},pdfauthor={Mohammed-Yassine Barnicha}}

\newtheorem{proposition}{Proposition}
\newtheorem{lemma}[proposition]{Lemma}
\theoremstyle{remark}
\newtheorem{remark}{Remark}

\newcommand{\E}{\mathbb{E}}
\newcommand{\Prob}{\mathbb{P}}
\newcommand{\dd}{\,\mathrm{d}}
\pgfplotsset{every axis/.append style={width=0.48\textwidth,height=5.2cm,grid=major,grid style={gray!25},scaled ticks=false,
  tick label style={font=\footnotesize},label style={font=\footnotesize},legend style={font=\scriptsize,draw=none,fill=white,fill opacity=0.8,text opacity=1}}}

\title{Floating-Strike Lookback Options\\[4pt]
  \large Closed Form, Exact-Bridge and Time-Stepping Monte Carlo, Greeks and a Code Review}
\author{Mohammed-Yassine Barnicha}
\date{October 2026}

\begin{document}
\maketitle

\begin{abstract}
\noindent We price European floating-strike lookback options under Black--Scholes with a continuous dividend, at issue and
during their life. Two Monte Carlo engines written in C++ are studied: one draws the terminal value and the running extremum
exactly, from the law of the Brownian bridge, in a single step; the other simulates paths on a time grid and corrects the
discrete-monitoring bias with the Broadie--Glasserman--Kou shift. We derive the law of the running extremum, the
Goldman--Sosin--Gatto closed form, and the scaling properties of the price: at issue the delta equals $V/S$, and the issue
price has no curvature, while a contract already held keeps a positive gamma $g''(1)/S$, a distinction the original Greeks
missed. We add an exact bridge on each time step, which removes the discretisation bias whatever the number of steps, quantify
the bias of each scheme, and list the defects found in the C++ code. Everything is reproduced by the interactive page
\url{https://yassinebarnicha.github.io/options-lookback/}, whose JavaScript engines match the C++ to $10^{-15}$.
\end{abstract}

\tableofcontents

%======================================================================
\section{Introduction}
A floating-strike lookback call pays $S_T-\min_{[0,T]}S$: the holder buys at the lowest price of the period. The put pays
$\max_{[0,T]}S-S_T$. Both payoffs are path-dependent but only through one running extremum, which makes them a good test bed
for Monte Carlo methods: the exact price is known in closed form, and the error of a discretely monitored simulation can be
measured exactly.

The project contains two independent C++ implementations, wrapped in a common DLL and called from Excel/VBA:
\begin{itemize}[nosep]
  \item \emph{engine A} (Barnicha) draws $(S_T,\text{extremum})$ exactly in one step by the Brownian bridge;
  \item \emph{engine B} (Chivet) simulates a time grid and applies the BGK correction to the discrete extremum.
\end{itemize}
This report gives the theory behind both (\cref{sec:law,sec:closed,sec:mc}), the scaling properties that determine the Greeks
(\cref{sec:homog}), a numerical study (\cref{sec:num}) and the code review (\cref{sec:review}).

%======================================================================
\section{Contract and model}
Under the risk-neutral measure, $S_t=S\exp(\mu t+\sigma W_t)$ with $\mu=b-\sigma^2/2$, where $b=r-q$ is the carry ($r$ rate, $q$
continuous dividend yield). Write $X_t=\ln(S_t/S)$, a Brownian motion with drift $\mu$ and volatility $\sigma$, and
$M_T=\max_{[0,T]}X$, $m_T=\min_{[0,T]}X$. For a contract issued in the past, let $M\ge S$ (put) or $m\le S$ (call) be the
extremum already observed. The payoffs at $T$ (time to maturity from now) are
\begin{equation}
  C_T=S_T-\min\big(m,\;S\,e^{m_T}\big),\qquad P_T=\max\big(M,\;S\,e^{M_T}\big)-S_T,
\end{equation}
with $m=M=S$ at issue. Prices are $V=e^{-rT}\E[\,\cdot\,]$; since $\E[e^{-rT}S_T]=Se^{-qT}$,
\begin{equation}\label{eq:decomp}
  V^{\mathrm{put}}=e^{-rT}\E\big[\max(M,Se^{M_T})\big]-S\,e^{-qT},\qquad V^{\mathrm{call}}=S\,e^{-qT}-e^{-rT}\E\big[\min(m,Se^{m_T})\big].
\end{equation}

%======================================================================
\section{The law of the running extremum}\label{sec:law}
\begin{proposition}[Joint law]\label{prop:joint}
For $y\ge\max(0,x)$,
\begin{equation}
  \Prob\big(M_T\ge y,\;X_T\in\mathrm dx\big)=\frac{1}{\sigma\sqrt T}\,\varphi\!\left(\frac{2y-x+\mu T}{\sigma\sqrt T}\right)e^{2\mu y/\sigma^2}\dd x .
\end{equation}
\end{proposition}
\begin{proof}
Without drift ($\mu=0$, $\sigma=1$), the reflection principle states that the path reflected after the first passage at $y$
is again a Brownian motion, so $\Prob(M_T\ge y,W_T\in\mathrm dx)=\Prob(W_T\in\mathrm d(2y-x))$. The drift is added by the
Girsanov density $\exp(\mu x/\sigma^2-\mu^2T/2\sigma^2)$, which depends on the end point only; combining the exponents gives
the formula.
\end{proof}
Integrating in $x$ gives the marginal law
\begin{equation}\label{eq:maxlaw}
  \Prob(M_T\le y)=\Phi\!\left(\frac{y-\mu T}{\sigma\sqrt T}\right)-e^{2\mu y/\sigma^2}\,\Phi\!\left(\frac{-y-\mu T}{\sigma\sqrt T}\right),\qquad y\ge0,
\end{equation}
and symmetrically $\Prob(m_T\ge y)=\Phi\big(\frac{-y+\mu T}{\sigma\sqrt T}\big)-e^{2\mu y/\sigma^2}\Phi\big(\frac{y+\mu T}{\sigma\sqrt T}\big)$ for $y\le0$.
Dividing the joint law by the density of $X_T$ removes the drift:
\begin{proposition}[Brownian bridge]\label{prop:bridge}
Conditionally on $X_0=x_0$ and $X_\tau=x_1$, for $y\ge\max(x_0,x_1)$,
\begin{equation}\label{eq:bridge}
  \Prob\Big(\max_{[0,\tau]}X\ge y\,\Big|\,x_0,x_1\Big)=\exp\!\left(-\frac{2(y-x_0)(y-x_1)}{\sigma^2\tau}\right),
\end{equation}
and the same holds for the minimum with $y\le\min(x_0,x_1)$.
\end{proposition}
The law does not depend on $\mu$: once both end points are known, the drift carries no information about the path.

%======================================================================
\section{The closed form}\label{sec:closed}
\begin{proposition}[Goldman, Sosin and Gatto; general form]\label{prop:gsg}
With $b_1=\dfrac{\ln(S/M)+(b+\sigma^2/2)T}{\sigma\sqrt T}$ and $b_2=b_1-\sigma\sqrt T$,
\begin{equation}
  V^{\mathrm{put}}=-S\,e^{(b-r)T}\Phi(-b_1)+M e^{-rT}\Phi(-b_2)
   +S\,e^{-rT}\frac{\sigma^2}{2b}\left[-\Big(\frac SM\Big)^{-2b/\sigma^2}\Phi\Big(b_1-\frac{2b}{\sigma}\sqrt T\Big)+e^{bT}\Phi(b_1)\right],
\end{equation}
and, with $a_1=\dfrac{\ln(S/m)+(b+\sigma^2/2)T}{\sigma\sqrt T}$, $a_2=a_1-\sigma\sqrt T$,
\begin{equation}
  V^{\mathrm{call}}=S\,e^{(b-r)T}\Phi(a_1)-m e^{-rT}\Phi(a_2)
   +S\,e^{-rT}\frac{\sigma^2}{2b}\left[\Big(\frac Sm\Big)^{-2b/\sigma^2}\Phi\Big(-a_1+\frac{2b}{\sigma}\sqrt T\Big)-e^{bT}\Phi(-a_1)\right].
\end{equation}
\end{proposition}
\begin{proof}[Sketch]
In \eqref{eq:decomp}, with $k=\ln(M/S)\ge0$,
$\E[\max(M,Se^{M_T})]=M\,\Prob(M_T\le k)+S\int_k^\infty e^y\,\mathrm dF(y)$, where $F$ is \eqref{eq:maxlaw}. The first term gives
$Me^{-rT}\Phi(-b_2)$ after discounting. In the integral, $e^y$ times the Gaussian part of $\mathrm dF$ is a Gaussian of shifted mean
(completing the square), which produces the $\Phi(b_1)$ terms, and $e^y$ times the reflected part, $e^{(1+2\mu/\sigma^2)y}$ times a
Gaussian, produces the power $(S/M)^{-2b/\sigma^2}$ and the factor $\sigma^2/2b$ through $1+2\mu/\sigma^2=2b/\sigma^2$.
\end{proof}
\begin{remark}[Zero carry]
Each bracket vanishes when $b\to0$ and the limit is finite. The general formula is evaluated at $b=\pm10^{-6}$ and averaged when
$|b|<10^{-6}$, an error of order $10^{-12}$; engine A codes the exact limit for $q=0$, $r\to0$; engine B divides by $2r$ and is
undefined at $r=0$.
\end{remark}
On the reference case ($S=100$, $r=5\,\%$, $q=0$, $\sigma=20\,\%$, $T=1$) the put is worth $14.290568$ and the call $17.216802$.

%======================================================================
\section{Homogeneity and the Greeks}\label{sec:homog}
\subsection{Scaling}
The dynamics is multiplicative and the payoffs are positively homogeneous in $(S,\text{extremum})$, hence
\begin{equation}
  V(\lambda S,\lambda m)=\lambda\,V(S,m)\quad(\lambda>0),\qquad V(S,m)=S\,g(m/S),
\end{equation}
for a function $g$ depending on $(r,q,\sigma,T)$. By Euler's theorem $V=S\,\partial_SV+m\,\partial_mV$.

\begin{lemma}\label{lem:neumann}
At $m=S$, $\partial_mV=0$ (one-sided derivative).
\end{lemma}
\begin{proof}
For the call, $\partial_m\min(m,Se^{m_T})=\mathbf 1\{m<Se^{m_T}\}$, so $\partial_mV=-e^{-rT}\Prob\big(\min_{[0,T]}S_t>m\big)$. As $m\uparrow S$
this probability tends to $\Prob(m_T>0)=0$: a Brownian motion goes below its starting point immediately. The put is symmetric.
\end{proof}

\begin{proposition}[Greeks at issue]\label{prop:greeks}
With $z=m/S$, $\Delta=\partial_SV=g(z)-z\,g'(z)$ and $\Gamma=\partial_{SS}V=z^2g''(z)/S$. At issue ($z=1$):
\begin{enumerate}[nosep,label=(\roman*)]
  \item $\Delta=g(1)=V/S$, the same for the issue price $S\mapsto V(S,S)$ and for the contract held (\cref{lem:neumann});
  \item the issue price $V(S,S)=S\,g(1)$ is linear in $S$: its second derivative is $0$;
  \item the contract held, whose extremum stays at $m$ when the spot moves, has $\Gamma=g''(1)/S>0$;
  \item vega, rho and theta are proportional to $S$.
\end{enumerate}
\end{proposition}
Point (iii) matters for hedging: right after issue the extremum is a fixed number in the contract, not a function of the spot.
On the reference case the held gamma is $0.0312$ for the put and $0.0455$ for the call, while the issue-price gamma is zero. Both C++ engines bump the spot \emph{and} the starting extremum, so they report the
gamma of (ii), not of (iii). The page computes (iii) by fixing the extremum (seasoned mode); \cref{fig:homog} shows both.

\begin{figure}[ht]\centering
\begin{tikzpicture}
\begin{axis}[xlabel={spot $S$},ylabel={$\Delta$},legend pos=south east,yticklabel style={/pgf/number format/fixed,/pgf/number format/precision=2}]
  \addplot[blue,very thick] table[x=S,y=delta]{data/homogene.dat}; \addlegendentry{$\Delta$ at issue}
  \addplot[orange,thick,dashed] table[x=S,y=VsurS]{data/homogene.dat}; \addlegendentry{$V/S$}
  \addplot[red,thick] table[x=S,y=deltavie]{data/homogene.dat}; \addlegendentry{$\Delta$, $M=110$ held}
\end{axis}
\end{tikzpicture}\hfill
\begin{tikzpicture}
\begin{axis}[xlabel={spot $S$},ylabel={$\Gamma$},legend pos=north east,yticklabel style={/pgf/number format/fixed,/pgf/number format/precision=3}]
  \addplot[blue,very thick] table[x=S,y=gamma]{data/homogene.dat}; \addlegendentry{issue price $V(S,S)$}
  \addplot[teal,thick] table[x=S,y=gamma_tenu]{data/homogene.dat}; \addlegendentry{held from issue, $g''(1)/S$}
  \addplot[red,thick] table[x=S,y=gammavie]{data/homogene.dat}; \addlegendentry{$M=110$ held}
\end{axis}
\end{tikzpicture}
\caption{Put, reference parameters. Left: at issue $\Delta=V/S=0.1429$ for every spot; with a maximum $M=110$ already observed,
$\Delta$ varies. Right: the issue price has zero gamma, the contract held does not.}\label{fig:homog}
\end{figure}

\subsection{Monte Carlo Greeks with common random numbers}
Both engines compute Greeks by bump and revalue on \emph{common random numbers} (CRN): the bumped prices reuse the same draws.
For a parameter $\theta$ and a payoff $Y(\theta,\omega)$ Lipschitz in $\theta$ with constant $L(\omega)$,
\[
  \operatorname{Var}\Big[\tfrac{Y(\theta+h)-Y(\theta-h)}{2h}\Big]\le\E[L^2]\quad\text{(CRN)},\qquad
  \operatorname{Var}\Big[\tfrac{Y(\theta+h)-Y'(\theta-h)}{2h}\Big]=\frac{\operatorname{Var}Y}{2h^2}\quad\text{(independent draws)},
\]
so CRN keeps the variance bounded as $h\to0$. At issue each path's payoff is itself homogeneous of degree one in $S$ (the
whole path and the starting extremum scale with $S$). Hence, path by path, the bumped delta equals $Y/S$ and the bumped gamma
is exactly zero: engine A returns $\Gamma=-1.1\cdot10^{-10}\pm1.4\cdot10^{-12}$ (rounding) and a delta whose standard error is
that of the price divided by $S$.

\paragraph{Theta conventions.} Engine A computes $(V(T-h)-V(T))/h$, a one-sided estimate of $-\partial_TV$; engine B computes
$+\partial_TV$ by central differences. The page keeps both conventions and offers a centred version of the first.

%======================================================================
\section{The Monte Carlo engines}\label{sec:mc}
\subsection{Engine A: exact bridge in one step}
Draw $X_T=\mu T+\sigma\sqrt T\,Z$, $Z\sim\mathcal N(0,1)$, then the extremum by inverting \eqref{eq:bridge} with $x_0=0$,
$x_1=X_T$ and a uniform $U$: solving $\exp(-2y(y-X_T)/\sigma^2T)=U$ for $y$,
\begin{equation}\label{eq:inverse}
  M_T=\tfrac12\Big(X_T+\sqrt{X_T^2-2\sigma^2T\ln U}\Big),\qquad m_T=\tfrac12\Big(X_T-\sqrt{X_T^2-2\sigma^2T\ln U'}\Big).
\end{equation}
The pair $(X_T,M_T)$ has exactly the law of \cref{prop:joint}: the estimator is unbiased with a single step. The antithetic
pair $(Z,-Z)$ with the same uniforms reduces the variance by a factor of about $2$ on the reference case (\cref{fig:conv}).
$U$ and $U'$ are drawn independently, which is correct for these payoffs (each uses one extremum) but would not give the joint
law of $(M_T,m_T)$ needed for a payoff such as $\max-\min$.

\subsection{Engine B: time stepping}
Paths are simulated on $n$ steps of length $\Delta t=T/n$ and the extremum is read on the grid. The discrete maximum is
smaller than the continuous one, and for a Brownian motion the gap is of order $\sqrt{\Delta t}$ (Asmussen, Glynn and Pitman
1995):
\begin{equation}
  \E\Big[\max_{[0,T]}X\Big]-\E\Big[\max_{k\le n}X_{k\Delta t}\Big]=\beta\,\sigma\sqrt{\Delta t}+o(\sqrt{\Delta t}),\qquad
  \beta=-\frac{\zeta(1/2)}{\sqrt{2\pi}}\approx0.5826.
\end{equation}
Broadie, Glasserman and Kou (1997, 1999) turn this into a correction: the discrete extremum is shifted,
$\max S\leftarrow e^{\beta\sigma\sqrt{\Delta t}}\max S$ and $\min S\leftarrow e^{-\beta\sigma\sqrt{\Delta t}}\min S$, which removes the
first-order term; the residual bias is $O(\Delta t)$. This is what the C++ does.

\subsection{Added: the exact bridge on each step}
Conditionally on the grid values, the pieces of path on different steps are independent Brownian bridges (Markov property).
The continuous maximum is the maximum over the steps of the bridge maxima, each of which can be drawn from \eqref{eq:bridge}
with its own uniform:
\begin{equation}
  \max_{[0,T]}X=\max_{k}\tfrac12\Big(X_k+X_{k+1}+\sqrt{(X_{k+1}-X_k)^2-2\sigma^2\Delta t\ln U_k}\Big).
\end{equation}
The scheme is then exact for continuous monitoring for any $n$; the grid only matters if the payoff needs intermediate values.
Its cost is one uniform and one square root per step.

%======================================================================
\section{Numerical study}\label{sec:num}
\subsection{Discretisation bias}
\Cref{fig:bias} shows the bias of engine B against the closed form, $20\,000$ paths per point, for the put. Without correction
the bias is $-1.25$ with $100$ steps (a $9\,\%$ underpricing) and decreases like $\sqrt{\Delta t}$. BGK overcorrects at coarse
grids ($+2.2$ with $2$ steps, $+0.35$ with $10$) and its residual decreases like $\Delta t$; it is within the statistical error from
about $50$ steps. The bridge per step has no visible bias at any $n$, including $n=2$.

\begin{figure}[ht]\centering
\begin{tikzpicture}
\begin{semilogxaxis}[width=0.8\textwidth,height=6.2cm,xlabel={time steps $n$},ylabel={price $-$ closed form},legend pos=south east,
  yticklabel style={/pgf/number format/fixed}]
  \addplot+[mark=*,thick,error bars/.cd,y dir=both,y explicit] table[x=n,y=brut,y error=brut_ic]{data/biais_put.dat}; \addlegendentry{no correction}
  \addplot+[mark=square*,thick,error bars/.cd,y dir=both,y explicit] table[x=n,y=bgk,y error=bgk_ic]{data/biais_put.dat}; \addlegendentry{BGK}
  \addplot+[mark=triangle*,thick,error bars/.cd,y dir=both,y explicit] table[x=n,y=pont,y error=pont_ic]{data/biais_put.dat}; \addlegendentry{exact bridge per step}
  \addplot[black,dashed,domain=2:500] {0};
\end{semilogxaxis}
\end{tikzpicture}
\caption{Discretisation bias of the time-stepping engine, put, reference parameters, with $95\,\%$ confidence intervals.}\label{fig:bias}
\end{figure}

\subsection{Convergence and antithetic variates}
\Cref{fig:conv} shows engine A against the closed form. Both estimators converge at the Monte Carlo rate $N^{-1/2}$; the
antithetic version has a standard error $1.42$ times smaller for the same number of normal draws (variance divided by $2.0$).

\begin{figure}[ht]\centering
\begin{tikzpicture}
\begin{semilogxaxis}[width=0.8\textwidth,height=6cm,xlabel={paths $N$},ylabel={put price},legend pos=north east,
  yticklabel style={/pgf/number format/fixed,/pgf/number format/precision=2}]
  \addplot+[mark=*,thick,error bars/.cd,y dir=both,y explicit] table[x=N,y=simple,y error=simple_ic]{data/convergence.dat}; \addlegendentry{plain}
  \addplot+[mark=square*,thick,error bars/.cd,y dir=both,y explicit] table[x=N,y=anti,y error=anti_ic]{data/convergence.dat}; \addlegendentry{antithetic}
  \addplot[black,dashed,domain=1000:200000] {14.290568}; \addlegendentry{closed form}
\end{semilogxaxis}
\end{tikzpicture}
\caption{Exact-bridge engine against the closed form, $95\,\%$ confidence intervals.}\label{fig:conv}
\end{figure}

\subsection{The sampler against the exact law}
\Cref{fig:law} compares the maxima drawn by \eqref{eq:inverse} ($200\,000$ paths) with the density from \eqref{eq:maxlaw}: the
histogram follows the exact law, including the mass near $0$ (paths that never rise much above the start).

\begin{figure}[ht]\centering
\begin{tikzpicture}
\begin{axis}[width=0.72\textwidth,height=5.4cm,xlabel={$y=\ln(\max S/S)$},ylabel={density},legend pos=north east,
  xticklabel style={/pgf/number format/fixed}]
  \addplot[ybar interval,fill=blue!25,draw=blue!50] table[x expr=\thisrow{y}-0.0083333,y=sim]{data/extremum.dat}; \addlegendentry{simulated}
  \addplot[red,very thick] table[x=y,y=exact]{data/extremum.dat}; \addlegendentry{exact}
\end{axis}
\end{tikzpicture}
\caption{Law of $M_T$ (reference parameters): bridge sampler against \eqref{eq:maxlaw}.}\label{fig:law}
\end{figure}

\subsection{Reference values}
\begin{table}[ht]\centering\small
\begin{tabular}{@{}lrrr@{}}\toprule
 & closed form & engine A & engine B, no correction\\
 & & ($10^4$ paths) & ($10^4$ paths $\times$ $100$ steps)\\\midrule
price & $14.2906$ & $14.2953\pm0.1357$ & $12.9073\pm0.1925$\\
$\Delta$ & $0.142906$ & $0.142953$ & $0.129073$\\
$\Gamma$ (issue price) & $0$ & $-1.1\cdot10^{-10}$ & $1.9\cdot10^{-11}$\\
vega & $87.170$ & $87.184$ & $79.460$\\
$\rho$ & $-49.379$ & $-49.371$ & $-48.900$\\\bottomrule
\end{tabular}
\caption{Put with the default settings of the Excel workbook ($S=100$, $r=5\,\%$, $\sigma=20\,\%$, $T=1$, seed 42), Monte Carlo
prices with $95\,\%$ confidence intervals. Engine B with BGK gives $14.2320$; with the bridge per step it is unbiased.}
\end{table}

%======================================================================
\section{Code review}\label{sec:review}
\begin{enumerate}
  \item \emph{The engine A program does not compile}: \texttt{Barnicha/app/main.cpp} has \texttt{using namespace lookback;} while the
        namespace is \texttt{Barnicha}. It is the executable called by the Excel macros.
  \item \emph{Engine B ignores the seed}: \texttt{getRandomSeed()} returns the clock, so two identical runs differ.
  \item \emph{Engine B, imperfect common random numbers}: \texttt{setRandomSeed} reseeds the generator but the normal law keeps the
        value held in reserve by the polar method, so the bumped prices do not use the same normals. This defeats the purpose of
        CRN and adds noise to every Greek.
  \item \emph{Engine B, standard error of $\rho$}: the estimate divides by $2hr$ (relative bump), the standard error by $2h$ only.
  \item \emph{Engine B, closed-form gamma four times too small}: divided by $(0.0002\,S)^2$ instead of $(0.0001\,S)^2$; invisible at
        issue, where the issue-price gamma is zero.
  \item \emph{Engine B, $r=0$}: the relative bump on $r$ and the closed form (division by $2r$) are undefined.
  \item \emph{Opposite theta conventions} between the two engines, one of them one-sided.
  \item \emph{Gamma at issue is not the hedging gamma}: both engines bump the starting extremum with the spot (\cref{prop:greeks}).
  \item \emph{Independent extremum draws} in engine A: fine for these payoffs, wrong for a payoff using both extrema.
  \item \emph{\texttt{out\_size} is never set} by the exported functions: the caller must know that there are 48 outputs.
\end{enumerate}
The JavaScript port fixes 2--6 (as options where the original behaviour is kept for comparison) and adds the dividend, the
seasoned extremum, the bridge per step, antithetic variates, a centred theta and adjustable bumps.

\section{Validation of the port}
The generator \texttt{std::mt19937} and libstdc++'s \texttt{generate\_canonical}, uniform and normal (polar Marsaglia)
distributions are reproduced bit for bit. With the same seed, engine A returns the prices of the C++ program compiled with g++
to $10^{-15}$ and engine B those of the C++ simulator (with and without BGK) to $10^{-12}$; the closed-form Greeks match the Excel
workbook to $10^{-10}$. With a dividend or a seasoned extremum, the Monte Carlo recovers the general closed form within its
confidence interval.

\section{Conclusion}
The exact bridge makes the lookback a one-step simulation: unbiased, cheap, and easily extended to any number of monitoring
steps by drawing one bridge extremum per step. Time stepping with BGK is a good approximation from about fifty steps but
overcorrects on coarse grids, while uncorrected discrete monitoring underprices by several percent. The scaling of the price
explains the Greeks: the delta at issue is $V/S$, and the gamma reported at issue is the curvature of the issue price, which is
zero, while the gamma of a contract held is $g''(1)/S$. The review found ten defects in the C++ code, the most important being
the uncompilable engine A program and the broken common random numbers of engine B.

\begin{thebibliography}{99}\small
\bibitem{agp} S.~Asmussen, P.~Glynn and J.~Pitman. Discretization error in simulation of one-dimensional reflecting Brownian
  motion. \emph{Annals of Applied Probability}, 5(4):875--896, 1995.
\bibitem{bgk97} M.~Broadie, P.~Glasserman and S.~Kou. A continuity correction for discrete barrier options.
  \emph{Mathematical Finance}, 7(4):325--349, 1997.
\bibitem{bgk99} M.~Broadie, P.~Glasserman and S.~Kou. Connecting discrete and continuous path-dependent options.
  \emph{Finance and Stochastics}, 3:55--82, 1999.
\bibitem{cv} P.~Glasserman. \emph{Monte Carlo Methods in Financial Engineering}. Springer, 2003.
\bibitem{gsg} M.~Goldman, H.~Sosin and M.~Gatto. Path dependent options: buy at the low, sell at the high.
  \emph{Journal of Finance}, 34(5):1111--1127, 1979.
\bibitem{haug} E.~G.~Haug. \emph{The Complete Guide to Option Pricing Formulas}. McGraw-Hill, 2nd ed., 2007.
\bibitem{shreve} S.~Shreve. \emph{Stochastic Calculus for Finance II}. Springer, 2004.
\end{thebibliography}
\end{document}
